\documentclass[sigplan,screen]{acmart}
\settopmatter{printfolios=true, printacmref=false}
\renewcommand\footnotetextcopyrightpermission[1]{} 
\makeatletter
\renewcommand\@formatdoi[1]{\ignorespaces}
\makeatother

\usepackage{setspace}
\usepackage{prelude}

% NOTE: Typographically improved footnotes
\let\origfootnote\footnote
\renewcommand{\footnote}[1]{\kern.06em\origfootnote{#1}}
\newcommand{\pfootnote}[1]{\kern-.06em\origfootnote{#1}}

\citestyle{acmauthoryear}

% WARN: Show citation keys; remove before submission
% \usepackage[notcite]{showkeys}
% \renewcommand*\showkeyslabelformat[1]{%
%   \fbox{\parbox[t]{\marginparwidth}{\smaller[3]\raggedleft\path{#1}}}}

% NOTE: Some typographical conventions
% If an argument for a function is not delimited, use \: to make the application more obvious
% For nested brackets, use \? or two, or even \! to reduce excessive spacings
% For superscripts like \ABS after a big closing semantic bracket, use \? to make the superscript closer
\begin{document}
\title[Retargeting an Abstract Interpreter for a New Language by Partial Evaluation]{Retargeting an Abstract Interpreter for a New Language by Partial Evaluation\\(Short Paper)}

\author{Jay Lee}
\orcid{0000-0002-2224-4861}
\email{jhlee@ropas.snu.ac.kr}
\affiliation{%
  \department{Department of Computer Science and Engineering}
  \institution{Seoul National University}
  \country{Korea}}

\author{Joongwon Ahn}
\orcid{0009-0007-0055-7297}
\email{jwahn@ropas.snu.ac.kr}
\affiliation{%
  \department{Department of Computer Science and Engineering}
  \institution{Seoul National University}
  \country{Korea}}

\author{Kwangkeun Yi}
\orcid{0009-0007-5027-2177}
\email{kwang@ropas.snu.ac.kr}
\affiliation{%
  \department{Department of Computer Science and Engineering}
  \institution{Seoul National University}
  \country{Korea}}

\copyrightyear{2026} 
\acmYear{2026}
\acmDOI{}
\setcopyright{rightsretained}
\acmConference[PEPM '26]{The 2026 ACM SIGPLAN International Workshop on Partial Evaluation and Program Manipulation}{January 13, 2026}{Rennes, France}

\keywords{Partial evaluation, Abstract interpretation}


\begin{abstract}
Implementing sound static analyzers for new languages requires significant time and effort.
We show that it is possible by partial evaluation to automatically retarget existing abstract interpreters for new languages.
Given an existing abstract interpreter for a source language and a definitional interpreter for a new target language, we can derive a correct abstract interpreter for the target language by partially evaluating the abstract interpreter with respect to the definitional interpreter.
We demonstrate our method on mini-languages and show that it is possible to mechanize the process with MetaOCaml.
Our approach suggests a promising direction to reduce the burden of building new analyzers from scratch.
\end{abstract}


\maketitle

\section{Introduction}\label{sec:prob}
Abstract interpretation \citep{CouCou77Abstract} provides a systematic framework to construct sound static analyzers by overapproximating a concrete semantics for a language.
While abstract interpreters can be designed in a principled manner following a recipe \citep{RivYi20Introduction}, it is time-consuming to model both a concrete semantics of a language and a corresponding abstract semantics.
This burden is compounded when an analysis designer repeats the process for every new target language.

We propose automatically deriving abstract interpreters for new languages from an existing abstract interpreter.
We show that partial evaluation \citep{Fut71Partial,Fut99Partial} of an abstract interpreter for a source language, using the semantics of target languages written in the source language in a form of a definitional interpreter \citep{Rey72Definitional}, leads to a retargeted sound abstract interpreter for the new target.
This eliminates the need to build analyzers for new languages from scratch.

Furthermore, we demonstrate using mini-languages that our approach is mechanizable by retargeting an abstract interpreter written in (BER)~MetaOCaml \citep{CalTahHuaLer03Implementing,Ole14Design}.


\section{Our Approach}\label{sec:overview}
There are three languages in play: a {\MET{}eta-language}, a {\SRC{}ource language}, and a new {\TGT{}arget language}.
The \ul{meta}-language~\MET{} is the implementation language of the abstract interpreter for language~\SRC, which also serves as the \ul{source} language of a definitional interpreter for the new analysis \ul{target}~\TGT.

Our goal is
\begin{description}
  \item[given] a concrete interpreter $\Intp ST$ for~\TGT{} written in~\SRC{} (a ``concrete \TGT*interpreter,'' see \cref{subsec:int}), and
  \item[given] an abstract interpreter $\Intp MS<\ABS>$ for~\SRC{} written in~\MET{} (an ``abstract \SRC*interpreter,'' see \cref{subsec:aint}), to
  \item[derive] an abstract interpreter $\Intp MT<\ABS>$ for~\TGT{} written in~\MET{} (an ``abstract \TGT*interpreter,'' see \cref{sec:correctness}).
\end{description}

A key observation is that partial evaluation~$\PE$ in the \MET*language (\cref{subsec:peval}) allows for a \emph{fully automatic} derivation of an abstract interpreter for~\TGT{}.
In \cref{sec:correctness}, we show that the retargeted abstract interpreter is indeed sound:
\OmitEnc
\OmitDec
\begin{restatable*}[Correctness of a Retargeted Abstract Interpreter]{thm}{correctness}\label{thm:correctness}
  Retargeting an abstract interpreter~$\Intp*MS<\ABS>$ with respect to an interpreter~$\Intp ST$ results in a sound abstract interpreter~$\Intp*MT<\ABS> \triangleq \Pe[\big]{\Intp MS<\ABS>}{\Enc[\big]SM{\Intp ST}}$.
\end{restatable*}

We finish by presenting a mechanized mini-example in \cref{sec:example} to illustrate our approach with MetaOCaml.

\section{The Ingredients}\label{sec:recipe}
\IncludeEnc
\IncludeDec
\begin{nota}
  We distinguish programs and values of each language using different typefaces and colors, e.g.,~$\mathM{p}$ for \MET-,~$\mathS{p}$ for \SRC-, and~$\mathT{p}$ for \TGT*programs.
  Explicit language name subscripts may be used when they are helpful, e.g.,~$\Dom MP$, $\Dom SP$, and~$\Dom TP$.
  Given domains~$\mathbb{D}_1$ and~$\mathbb{D}_2$, we denote the product as~$\mathbb{D}_1 \times \mathbb{D}_2$, disjoint union as~$\mathbb{D}_1 + \mathbb{D}_2$, and (Scott-)continuous function domain as~$\mathbb{D}_1 \domto \mathbb{D}_2$, where~$\times$ has the highest precedence, $+$ has lower precedence, and~$\domto$ the lowest; $\times$ and~$+$ are left-associative, whereas~$\domto$ is right-associative~\citep{Rey98Theories}.
  We denote the set of partial functions from~$X$ to~$Y$ as~$X \parto Y$.
  Angle brackets denote semantic tuples, e.g., $\Tup{v_1, v_2} \in \mathbb{D}_1 \times \mathbb{D}_2$ when $v_1 \in \mathbb{D}_1$ and $v_2 \in \mathbb{D}_2$.
  $\scond{\cdot}{\cdot}{\cdot}$ is the semantic conditional operator.
\end{nota}

We now present the ingredients of our recipe.
While our recipe is tuned to be readily applied to the mini-example of \cref{sec:example}, exact signatures for semantic functions and domains may be adapted to other forms.
\begin{defn}[Syntax and Semantics]\label{def:synsem}
  Let~\(\Dom MV\) be the set of values in~\MET,
  where \(\Dom MV\) is closed under product $\Dom MV \times \Dom MV \subseteq \Dom MV$ and continuous functions $(\Dom MV \domto \Dom MV) \subseteq \Dom MV$;
  and~\(\Dom SS\) and~\(\Dom TS\) be the sets of states in~\SRC{} and~\TGT.
  Let~\(\Dom MP\),~\(\Dom SP\), and~\(\Dom TP\) be programs of each language~\MET,~\SRC, and~\TGT.
  Then we have concrete semantics functions
  \begin{equation*}
    \Sem M{p} \in \Dom MV, \qquad
    \Sem S{p} \in \Dom SS \domto \Dom SS,\text{ and } \quad
    \Sem T{p} \in \Dom TS \domto \Dom TS
  \end{equation*}
  for all ~\(\mathM p \in \Dom MP\),~\(\mathS p \in \Dom SP\), and~\(\mathT p \in \Dom TP\).
\end{defn}
Moreover, each domain of defining-language values should be able to encode the programs and the values of the defined-language---this is expressed in \cref{def:encdec}.


\subsection{An Interpreter for the New Language}\label{subsec:int}
The first ingredient is the semantics of the new language~\TGT{}: a definitional interpreter~$\Intp ST$.
This is written in the source language~\SRC{} for which we have an abstract interpreter~$\Intp*MS<\ABS>$.

For an interpreter to accept programs of the target language, we need to embed the program and the initial state of the target language into the source language.
\begin{defn}[Encoders and Decoders]\label{def:encdec}
  An \emph{input encoder} $\Enc TS{\cdot, \cdot} \in \Dom TP \times \Dom TS \domto \Dom SS$ is a function that encodes a \TGT*program and a \TGT*state into an \SRC*state.
  This induces an \emph{output decoder}~$\Dec ST{} \in \Dom SS \parto \Dom TS$, which is a partial function, where
  \begin{equation*}
    \Dec[\big]ST{\Enc TS{p,\Stt T}} = \Stt T
  \end{equation*}
  for all $\mathT{p} \in \Dom TP$ and $\Stt T \in \Dom TS$.
  We overload the notation for the lifted decoder, that is, for any \(\mathS{\overline{\Stt S}} \in \Pow{\Dom SS}\) we have
  \begin{equation*}
    \Dec ST{\overline{\Stt S}} = \{ \Dec ST{\Stt S} \mid \Stt S \in \mathS{\overline{\Stt S}} \}. \qedhere
  \end{equation*}
\end{defn}
In practice, encoders and decoders are usually part of an interpreter as parsers, and we can assume that they are available.
For brevity, we shall omit the subscripts and superscripts of encoders and decoders as they can be unambiguously inferred.

%% NOTE: Omitting encoders and decoders for brevity, but keep them for correct typesetting
%\OmitEnc
%\OmitDec
\OmitSubEnc
\OmitSubDec

\begin{defn}[Concrete Interpreter]\label{def:int}
  A concrete \TGT*interpreter~$\Intp ST \in \Dom SP$ is an \SRC*program that accepts an \SRC*state that encodes a \TGT*program and a \TGT*state, and returns an \SRC*value that encodes the evaluation result.
  That is, $\Intp ST$ satisfies
  \begin{equation}\label{eq:int}
    %\Sem[\big]S{\Intp ST} \mathS{\Tup{\Enc TS{e}, \Enc TS{i}}} = \Enc[\big]TS{\Sem T{e}\: i}
    \Dec[\big]ST{\Sem[\big]S{\Intp ST} \Enc TS{p,\Stt T}} = \Sem T{p}\: \Stt T
  \end{equation}
  for all $\mathT{p} \in \Dom TP$ and $\Stt T \in \Dom TS$.
\end{defn}

Note that \cref{eq:int} enforces a few requirements on~\SRC.
\SRC*language should be expressive enough to encode the programs and values of \TGT,
and to write an interpreter for \TGT.
%Note that \cref{eq:int} enforces a few requirements on the expressiveness of the domains.
%\SRC*language should to be able to construct and represent a tuple, as well as to encode the programs and values of \TGT.
%And of course, \SRC*language should be able to write an interpreter for \TGT.
%These explain the requirements in \cref{def:synsem}.


\subsection{An Existing Abstract Interpreter}\label{subsec:aint}
Now we need the second ingredient: an existing abstract interpreter for~\SRC,~$\Intp*MS<\ABS>$.
This is written in the meta-language~\MET{} for which we have a partial evaluator~$\PE$.

We define the abstract domain and the concretization function for the source language~\SRC{} in which the abstract interpreter~$\Intp*MS<\ABS>$ operates.
\begin{defn}[Abstract Domain and Concretization]\label{def:conc}
  The abstract domain~$\Dom*SS<\ABS>$ of~\SRC{} approximates the concrete domain~$\Dom*SS$ with an approximation partial order~$\leS$.
  An abstract state can be concretized to a set of concrete states using a monotonic concretization function~$\Conc S \in \Dom*SS<\ABS> \domto \Pow{\Dom SS}$, where for any~$\smash{\mathS{{\Stt S}_1^\ABS}}, \smash{\mathS{{\Stt S}_2^\ABS}} \in \Dom*SS<\ABS>$, we have \[\mathS{{\Stt S}_1^\ABS} \leS \mathS{{\Stt S}_2^\ABS} \implies \Conc S \: \mathS{{\Stt S}_1^\ABS} \subseteq \Conc S \: \mathS{{\Stt S}_2^\ABS}.\]
  Moreover, we overload the notation of encoder and decoder for programs and abstract states, that is, we have \(\Enc SM{\cdot} \in \Dom SP + \Dom*SS<\ABS> \domto \Dom MV\), \(\Enc SM{\cdot,\cdot} \in \Dom SP \times \Dom*SS<\ABS> \domto \Dom MV \times \Dom MV\) and \(\Dec MS{} \in \Dom MV \parto \Dom*SS<\ABS>\), where
  \begin{equation*}
    \Enc SM{p,\Stt S<\ABS>} = \Tup{\Enc SM{p}, \Enc SM{\Stt S<\ABS>}}
  \end{equation*}
  and
  \begin{equation*}
    \Dec[\big]MS{\Enc SM{\Stt S<\ABS>}} = \Stt S<\ABS>. \qedhere
  \end{equation*}
\end{defn}

With \cref{def:conc}, we can define a correctness condition for an abstract \SRC*interpreter~$\Intp MS<\ABS>$.
\begin{defn}[Abstract Interpreter]\label{def:aint}
  An abstract \SRC*interpreter~$\Intp*MS<\ABS> \in \Dom MP$ is an \MET*program that accepts an \SRC*program and an \SRC*state, and returns a sound approximation of the evaluation result with respect to the concretization~$\Conc S$, as defined in \cref{def:conc}.
  That is, $\Intp*MS<\ABS>$~satisfies
  \begin{equation}\label{eq:aint}
    \Stt S \in \Conc S \: \Stt S<\ABS>
    \implies
    \Sem S{p}\: \Stt S \in \Conc S \: \Dec[\big]MS{\Sem[\big]M{\Intp MS<\ABS>} \Enc SM{p,\Stt S<\ABS>}}
  \end{equation}
  for all~$\mathS{p} \in \Dom SP$,~$\Stt*S<\ABS> \in \Dom*SS<\ABS>$ and~$\Stt S \in \Dom*SS$.
\end{defn}

Again, \cref{eq:aint} imposes requirements on the expressiveness of \MET.
The meta-language~\MET{} should be expressive enough to encode the programs and values of \SRC, and write an abstract interpreter for \SRC.
In addition,~\MET{} should be able to construct and represent tuples and functions, which explains the requirements in \cref{def:synsem}.


\subsection{Retargeting using a Partial Evaluator}\label{subsec:peval}
The final ingredient for automatically deriving a retargeted abstract interpreter is a partial evaluator for \MET: $\PE \in \Dom MP \times \Dom MV \domto \Dom MP$.
\begin{defn}[Partial Evaluator]\label{def:peval}
  A partial evaluator~$\PE \in \Dom MP \times \Dom MV \domto \Dom MP$ accepts two inputs,
  \begin{enumerate}
    \item a program~$\mathM{p} \in \Dom MP$ that accepts~$\mathM{\Tup{v_1, v_2}} \in \Dom MV$, and
    \item a value~$\mathM{v_1} \in \Dom MV$,
  \end{enumerate}
  and returns a specialized program that accepts the remaining value~$\mathM{v_2} \in \Dom MV$.
  That is, $\PE$~satisfies
  \begin{equation}\label{eq:peval}
    \Sem M{\Pe[\big]{p}{v_1}}\: \mathM{v_2} = \Sem M{p} \mathM{\Tup{v_1, v_2}}
  \end{equation}
  for all~$\mathM{p} \in \Dom MP$ and~$\mathM{v_1}, \mathM{v_2} \in \Dom MV$.
\end{defn}

From \cref{eq:peval}, we can specialize an abstract interpreter with a target program by substituting~$\Intp MS<\ABS>$ for~$\mathM{p}$:
\begin{equation}\label{eq:fut}
  \Sem[\big]M{\Pe[\big]{\Intp MS<\ABS>}{\Enc SM{p}}}\: \Enc SM{\Stt S<\ABS>}
  = \Sem[\big]M{\Intp MS<\ABS>}\: \Enc SM{p, \Stt S<\ABS>}.
\end{equation}
This is essentially the first Futamura projection extended to abstract interpreters.


\section{The Correctness Theorem}\label{sec:correctness}
Before we show the \hyperref[thm:correctness]{main theorem}, we prove the correctness of a meta-level analysis.

\begin{lem}[Correctness of a Meta-level Analysis]\label{lem:meta}
  Analyzing a \TGT*program using an abstract interpreter~$\Intp MS<\ABS>$ along with the concrete interpreter~$\Intp ST$ is sound.
\end{lem}
\begin{proof}
  Suppose we are given a \TGT*program~$\mathT{p}$ and a \TGT*state~$\Stt T$.
  Then by substituting $\Intp ST$ for $\mathS{p}$ and $\mathS{{\Enc TS{p,\Stt T}}}$ for $\Stt S$ in \cref{eq:aint},
  \begin{align*}
    \Enc TS{p,\Stt T} \in \Conc S \: \Stt S<\ABS>
    & \implies
    \Sem[\big]S{\Intp ST} \Enc TS{p,\Stt T}
    \in
    \Conc S\: \Dec[\big]MS{\Sem[\big]M{\Intp MS<\ABS>} \: \Enc SM{\Intp ST, \Stt S<\ABS>}} \\
    & \implies
    \Dec[\big]ST{\Sem[\big]S{\Intp ST} \Enc TS{p,\Stt T}}
    \in
    \Dec[\big]ST{\Conc S\: \Dec[\big]MS{\Sem[\big]M{\Intp MS<\ABS>} \: \Enc SM{\Intp ST, \Stt S<\ABS>}}}
  \end{align*}
  Then from \cref{eq:int}, $\Dec[\big]ST{\Sem[\big]S{\Intp ST} \Enc TS{p,\Stt T}}$ is precisely $\Sem T{p}\: \Stt T$.
  Thus the analysis~$\Sem[\big]M{\Intp MS<\ABS>} \: \Enc SM{\Intp ST, \Stt S<\ABS>}$ subsumes the concrete evaluation $\Sem T{p}\: \Stt T$ when decoded.
\end{proof}

\correctness
\begin{proof}
  Given a partial evaluator~$\PE$ that satisfies \cref{eq:peval}, we can partially evaluate the abstract interpreter~$\Intp*MS<\ABS>$ with respect to the interpreter~$\Intp*ST$ as
  \begin{equation*}
    \Intp*MT<\ABS> \triangleq \Pe[\big]{\Intp*MS<\ABS>}{\Enc[\big]SM{\Intp*ST}}.
  \end{equation*}

  We now show that the specialized~$\Intp*MT<\ABS>$ is indeed a correct static analyzer of~\TGT.
  Fix a \TGT*program~$\mathT{p}$ and a \TGT*state~$\Stt T$.
  Given \(\Stt S<\ABS>\) such that \(\Enc TS{p,\Stt T} \in \Conc S \: \Stt S<\ABS>\),
  from \cref{eq:fut},
  \begin{equation*}
    \Sem[\big]M{\Intp*MT<\ABS>}\: \Enc SM{\Stt S<\ABS>} = \Sem[\big]M{\Pe[\big]{\Intp*MS<\ABS>}{\Enc SM{\Intp*ST}}}\: \Enc SM{\Stt S<\ABS>} = \Sem[\big]M{\Intp*MS<\ABS>}\: \Enc SM{\Intp*ST, \Stt S<\ABS>}
  \end{equation*}
  and from \cref{lem:meta}, this soundly approximates the evaluation
  \begin{equation*}
    \Sem T{p} \: \Stt T
    \in
    \Dec[\big]ST{\Conc S\: \Dec[\big]MS{\Sem[\big]M{\Intp MS<\ABS>} \: \Enc SM{\Intp ST, \Stt S<\ABS>}}}
  \end{equation*}
  Thus the specialized abstract interpreter~$\Intp*MT<\ABS>$ is correct.
\end{proof}


\section{A Mechanized Example}\label{sec:example}
We now present a concrete example that demonstrates how our approach retargets an abstract interpreter~$\Intp*MS<\ABS>$ for a source language~\SRC{} to an abstract interpreter~$\Intp*MT<\ABS>$ for a target language~\TGT.
For mechanization, we use MetaOCaml to be the meta-language~\MET.
This example illustrates how our method works in practice and highlights the properties of the resulting retargeted abstract interpreter.

Consider mini-languages~\TGT{} and~\SRC:
\begin{center}
  \begin{bnf}
    \mathT{p} : $\Dom TP$ ::= \Tadd{n} // \Tmult{n} // \Tseq{p}{p}
  \end{bnf}
\end{center}
\begin{center}
  \begin{bnf}
    \mathS{e} : $\Dom SE$ ::=
    | $\mathS{n}$ // \Sread{e} // \Sadd{e}{e} // \Smult{e}{e} // \Sequ{e}{e}
    ;;
    \mathS{p} : $\Dom SP$ ::=
    | \Sskip{} // \Sasgn{e}{e} // \Sseq{p}{p}
    | \Sif{e}{p}{p} // \Swhile{e}{p}
  \end{bnf}
\end{center}
\TGT{}~is an assembly-like language with a single register, e.g., \Tseq{\Tadd{42}}{\Tmult{3}} returns~\((\Stt T + 42) \times 3\), given an initial register value~\(\Stt T\).
%\TGT{}~is a set of single instructions that either add or multiply a number to a user input, e.g., \Tadd{42} adds 42 to an input and \Tmult{42} multiplies an input by 42.
\SRC{}~is a structured language with a memory that allows arbitrary assignment at natural number addresses, e.g, \Sasgn{\textS{(\Sadd{13}{37})}}{0} transforms the value at address~50 to~0.

% \IncludeEnc
% \IncludeDec
Given the integer value domains $\Dom TV = \Dom SV = \mathbb{Z}$ and state domains $\Dom TS = \Dom TV$ as a single register value and $\Dom SS = \mathS{\Dom S{\overline V}}$ as an array of memory values,
the semantics~$\Sem T{}$ and~$\Sem S{}$ and the encoder~$\Enc TS{\cdot, \cdot}$ are given as follows:
\begin{mathpar}
  \Sem T{p} \in \Dom TV \domto \Dom TV \and
  \Enc TS{p, \Stt T} \in \Dom S{\overline V} \\

  \Sem T{\Tadd{n}} \triangleq \func{\Stt T}{\mathT{n} + \Stt T} \and
  \Sem T{\Tmult{n}} \triangleq \func{\Stt T}{\mathT{n} \times \Stt T} \and
  \Sem T{\Tseq{p_1}{p_2}} \triangleq \Sem T{p_2} \Comp \Sem T{p_1} \and
  \Enc TS{\Tseq{p_1}{\Tseq{...}{p_k}}, \Stt T} \triangleq \Stt T \cons 2 \cons \Concat_{\mathT{i}=\mathT{1}}^{\mathT{k}} \bigl[\scond{\mathT{p_i} = \Tadd{n}}{1}{2}, \mathT{n}\bigr] \concat [0] \\

  \Sem S{e} \in \Dom SS \domto \Dom SV \and
  \Sem S{p} \in \Dom SS \domto \Dom SS \\

  \Sem S{n} \triangleq \func{\Stt S}{\mathS{n}} \and
  \Sem S{\Sread{e}} \triangleq \func{\Stt S}{\Stt S[\Sem S{e} \Stt S]} \and
  \Sem S{\Sadd{e_1}{e_2}} \triangleq \func{\Stt S}{\Sem S{e_1} \Stt S + \Sem S{e_2} \Stt S} \and
  \Sem S{\Smult{e_1}{e_2}} \triangleq \func{\Stt S}{\Sem S{e_1} \Stt S \times \Sem S{e_2} \Stt S} \and
  \Sem S{\Sequ{e_1}{e_2}} \triangleq \func{\Stt S}{\scond{\Sem S{e_1} \Stt S = \Sem S{e_2} \Stt S}{1}{0}} \\

  \Sem S{\Sskip{}} \triangleq \func{\Stt S}{\Stt S} \and
  \Sem S{\Sasgn{e_1}{e_2}} \triangleq \func{\Stt S}{\Stt S[\Sem S{e_1} \mapsto \Sem S{e_2}]} \and
  \Sem S{\Sseq{p_1}{p_2}} \triangleq \Sem S{p_2} \Comp \Sem S{p_1} \and
  \Sem S{\Sif{e}{p_1}{p_2}} \triangleq \func{\Stt S}{\scond{\Sem S{e} \Stt S \ne 0}{\Sem S{p_1} \Stt S}{\Sem S{p_2} \Stt S}} \and
  \Sem S{\Swhile{e}{p}} \triangleq \func{\Stt S}{\scond{\Sem S{e} \Stt S = 0}{\Stt S}{\Sem S{\Swhile{e}{p}} \Comp \Sem S{p} \: \Stt S}}
\end{mathpar}

To elaborate on the encoding~$\Enc TS{p, \Stt T}$, we place the register value~$\Stt T$ and the program counter---initially~2---at the first and second addresses.
Each instruction is stored with an alignment of two words: the first cell stores the operation code---\textT{ADD} is~1 and \textT{MUL} is~2---and the second stores the argument.
The program finishes with the operation code~0 for halting.
We omit the presentation of~$\Dom MP$,~$\Dom MV$,~$\Enc SM{\cdot}$, and~$\Sem M{e}$ of MetaOCaml~(\MET) for brevity.
The encoder~$\Enc SM{\cdot}$ is a typical parser that embeds the AST of an \SRC*program or an \SRC*value into its corresponding \MET*datatype representation.

Now we can define the concrete \TGT*interpreter
\begin{srccode}
$\Intp ST ={}$while (*PC) {
  if (*PC == 1) { R = R + *(PC + 1) }
  else if (*PC == 2) { R = R * *(PC + 1) }
  else { skip };
  PC = PC + 2 }
\end{srccode}
where macros $\mathS{R} = \Sread{0}$ and $\mathS{PC} = \Sread{1}$ are used for convenience.

The abstract interpreter~$\Intp*MS<\ABS>$ parameterized by the base value abstraction is shown below in (pseudo-)MetaOCaml. %in \cref{fig:impl}.
Brackets \textsf{.<\textellipsis>.} wraps code to generate, and an escape \textsf{.\textasciitilde} represents a hole inside the code template.
% \begin{figure}

 \begin{metcode}
let $\mathM{\Intp*MS<\ABS>(\textS{p}, \Stt S<\ABS>)}$ = $\mathM{\textM{evalp}^\ABS(\textS{p}, \Stt S<\ABS>)}$
where $\mathM{\textM{evale}^\ABS(\mathS{e}, \Stt S<\ABS>)}$ = match $\mathS{e}$ with
  | $\mathS{n} \to$ .<$\eta(\mathS{n})$>.
  | $\Sread{e} \to$ .< .~$\Stt S<\ABS>$[.~$\mathM{\textM{evale}^\ABS(\mathS{e}, \Stt S<\ABS>)}$] >.
  | $\Sadd{e_1}{e_2} \to$ .< .~$\mathM{\textM{evale}^\ABS(\mathS{e_1}, \Stt S<\ABS>)}$ $\mathbin{+^\ABS}$ .~$\mathM{\textM{evale}^\ABS(\mathS{e_2}, \Stt S<\ABS>)}$ >.
  | $\Smult{e_1}{e_2} \to$ .< .~$\mathM{\textM{evale}^\ABS(\mathS{e_1}, \Stt S<\ABS>)}$ $\mathbin{\times^\ABS}$ .~$\mathM{\textM{evale}^\ABS(\mathS{e_2}, \Stt S<\ABS>)}$ >.
  | $\Sequ{e_1}{e_2} \to$ .< .~$\mathM{\textM{evale}^\ABS(\mathS{e_1}, \Stt S<\ABS>)}$ $\mathbin{=^\ABS}$ .~$\mathM{\textM{evale}^\ABS(\mathS{e_2}, \Stt S<\ABS>)}$ >.
and $\mathM{\textM{evalp}^\ABS(\mathS{p}, \Stt S<\ABS>)}$ = match $\mathS{p}$ with
  | $\Sskip{} \to$ $\Stt S<\ABS>$
  | $\Sasgn{e_1}{e_2} \to$ .< .~$\Stt S<\ABS>$[ .~$\mathM{\textM{evale}^\ABS(\mathS{e_1}, \Stt S<\ABS>)}$ $\mapsto^\ABS$ .~$\mathM{\textM{evale}^\ABS(\mathS{e_2}, \Stt S<\ABS>)}$ ] >.
  | $\Sseq{p_1}{p_2} \to$ .< .~$\mathM{\textM{evalp}^\ABS(\mathS{p_2},
             \textM{evalp}^\ABS(\mathS{p_1}, \Stt S<\ABS>))}$ >.
  | $\Sif{e}{p_t}{p_f} \to$
    .< let $\mathS{v^\ABS}$ = .~$\mathM{\textM{evale}^\ABS(\mathS{e}, \Stt S<\ABS>)}$ in
       $\mathcal{F}^\ABS_{\ne 0}(\mathS{v^\ABS}, \textM{.\textasciitilde evalp}^\ABS(\mathS{p_t}, \Stt S<\ABS>))
          \mathbin{\sqcup^\ABS}
          \mathcal{F}^\ABS_{=0}(\mathS{v^\ABS}, \textM{.\textasciitilde evalp}^\ABS(\mathS{p_f}, \Stt S<\ABS>))$ >.
  | $\Swhile{e}{p_b} \to$
    .< let $\mathS{v^\ABS}$ = .~$\mathM{\textM{evale}^\ABS(\mathS{e}, \Stt S<\ABS>)}$ in
      .~$\mathM{\textM{evalp}^\ABS(\mathS{p},
                \textM{evalp}^\ABS(\mathS{p_b},
                   \textM{.< } \mathcal{F}^\ABS_{\ne 0}(\mathS{v^\ABS}, \Stt S<\ABS>)\textM{ >.}))} \mathbin{\sqcup^\ABS} \mathcal{F}^\ABS_{=0}(\mathS{v^\ABS}, \Stt S<\ABS>)$ >.
 \end{metcode}

% \begin{metcode}
% let $\mathM{\Intp*MS<\ABS>(\textS{p}, \Stt S<\ABS>)}$ = $\mathM{\textM{evalp}^\ABS(\textS{p}, \Stt S<\ABS>)}$
% where $\mathM{\textM{evale}^\ABS(\mathS{e}, \Stt S<\ABS>)}$ = match $\mathS{e}$ with
%   | $\mathS{n} \to \eta(\mathS{n})$
%   | $\Sread{e'} \to \Stt S<\ABS>[\mathM{evale}^\ABS(\mathS{e'},\Stt S<\ABS>)]^\ABS$
%   | $\Sadd{e_1}{e_2} \to$
%       .<.~$\mathM{\textM{evale}^\ABS(\mathS{e_1}, \Stt S<\ABS>)}$ $\mathbin{+^\ABS}$ .~$\mathM{\textM{evale}^\ABS(\mathS{e_2}, \Stt S<\ABS>)}$>.
%   | $\Smult{e_1}{e_2} \to
%       \mathM{\textM{evale}^\ABS(\mathS{e_1}, \Stt S<\ABS>) \mathbin{\times^\ABS}
%              \textM{evale}^\ABS(\mathS{e_2}, \Stt S<\ABS>)}$
%   | $\Sequ{e_1}{e_2} \to
%       \mathM{\textM{evale}^\ABS(\mathS{e_1}, \Stt S<\ABS>) \mathbin{=^\ABS}
%              \textM{evale}^\ABS(\mathS{e_2}, \Stt S<\ABS>)}$
% and $\mathM{\textM{evalp}^\ABS(\mathS{p}, \Stt S<\ABS>)}$ = match $\mathS{p}$ with
%   | $\Sskip{} \to \Stt S<\ABS>$
%   | $\Sasgn{e_1}{e_2} \to$
%       $\Stt S<\ABS>[\textM{evale}^\ABS(\mathS{e_1}, \Stt S<\ABS>)
%       \mapsto^\ABS \textM{evale}^\ABS(\mathS{e_2}, \Stt S<\ABS>)]$
%   | $\Sseq{p_1}{p_2} \to$
%       $\mathM{\textM{evalp}^\ABS(\mathS{p_2}, 
%       \textM{evalp}^\ABS(\mathS{p_1}, \Stt S<\ABS>))}$
%   | $\Sif{e}{p_t}{p_f} \to$
%       let $\mathS{v^\ABS} = \textM{evale}^\ABS(\mathS{e}, \Stt S<\ABS>)$ in
%       $\mathM{\mathcal{F}^\ABS_{\ne 0}(\mathS{v^\ABS}, \textM{evalp}^\ABS(\mathS{p_t}, \Stt S<\ABS>))
%        \mathbin{\sqcup^\ABS}
%        \mathcal{F}^\ABS_{= 0}(\mathS{v^\ABS}, \textM{evalp}^\ABS(\mathS{p_f}, \Stt S<\ABS>))}$
%   | $\Swhile{e}{p_b} \to$
%       let $\mathS{v^\ABS}$ = $\textM{evale}^\ABS(\mathS{e}, \Stt S<\ABS>)$ in
%       $\textM{evalp}^\ABS(\mathS{p},
%         \textM{evalp}^\ABS(\mathS{p_b}, \mathcal{F}^\ABS_{\ne 0}(\mathS{v^\ABS}, \Stt S<\ABS>)))
%         \mathbin{\sqcup^\ABS}
%         \mathcal{F}^\ABS_{= 0}(\mathS{v^\ABS}, \Stt S<\ABS>)$
% \end{metcode}
% \caption{A (pseudo) MetaOCaml implementation of an abstract interpreter~$\Intp*MS<\ABS>$.}
% \label{fig:impl}
% \end{figure}
The parametrizable knob is given by the extraction function $\eta$ \citep{NieNieHan99Principles,DarHor19Constructive} that extracts the abstract value from a single concrete value.
Abstract operators are defined correspondingly: abstract operators~$+^\ABS$,~$\times^\ABS$,~$=^\ABS$ correspond to~$+$,~$\times$,~$=$, respectively; $\sqcup^\ABS$~is an abstract join; and~$\mathcal{F}^\ABS_{\ne 0}$ and~$\mathcal{F}^\ABS_{= 0}$ filter abstract values based on the predicate~$\ne 0$ and~$= 0$, respectively \citep{RivYi20Introduction}.
% We preserve the tuple structure of the values and use the concrete projections~$\pi_1$ and~$\pi_2$ and concrete tuple constructor~$\mathM{(\cdot, \cdot)}$.

To ensure termination, we have also modified the above code into an open-recursive style in the actual implementation, \textit{\`a la} Abstracting Definitional Interpreters \citep{DarLabNguVan17Abstracting}.

Performing the specialization by applying \(\Intp*ST\) to the above function, we get the retargeted abstract interpreter~$\Intp MT<\ABS> = \Pe[\big]{\Intp MS<\ABS>}{\Enc[\big]SM{\Intp ST}}$ as\footnote{The open-recursive style has been manually restored to a recursive style for presentation.}
%\(\func{\Stt S<\ABS>}{\Intp*MS<\ABS>(\Intp*ST, \textM{.< }\Stt S<\ABS> \textM{ >.})}\), 
\begin{metcode}
let $\mathM{{\Intp MS<\ABS>}_{\Intp ST}(\Stt S<\ABS>)}$ =
  let $\mathS{v_1} = \mathM{\Stt S<\ABS>[\Stt S<\ABS>[\eta(1)]]}$ in
  let $\mathS{\Stt S^\ABS_2} = \mathM{\mathcal{F}^\ABS_{\ne 0}(\mathS{v_1}, \Stt S<\ABS>)}$ in
  let $\mathS{v_3} = \mathM{\mathS{\Stt S^\ABS_2}[\mathS{\Stt S^\ABS_2}[\eta(1)]] =^\ABS \eta(0)}$ in
  let $\mathS{\Stt S^\ABS_4} = \mathM{\mathcal{F}^\ABS_{\ne 0}(\mathS{v_3}, \mathS{\Stt S^\ABS_2})}$ in let $\mathS{\Stt S^\ABS_5} = \mathM{\mathcal{F}^\ABS_{=0}(\mathS{v_3}, \mathS{\Stt S^\ABS_2})}$ in
  let $\mathS{v_6} = \mathM{\mathS{\Stt S^\ABS_5}[\mathS{\Stt S^\ABS_5}[\eta(1)]] =^\ABS \eta(2)}$ in
  let $\mathS{\Stt S^\ABS_7} = \mathM{\mathcal{F}^\ABS_{\ne 0}(\mathS{v_6}, \mathS{\Stt S^\ABS_5})}$ in let $\mathS{\Stt S^\ABS_8} = \mathM{\mathcal{F}^\ABS_{=0}(\mathS{v_6}, \mathS{\Stt S^\ABS_5})}$ in
  let $\mathS{\Stt S^\ABS_9} =$
    $\mathM{\mathS{\Stt S^\ABS_4}\bigl[\eta(0) \mapsto^\ABS \mathS{\Stt S^\ABS_4}[\eta(0)] +^\ABS \mathS{\Stt S^\ABS_4}[\mathS{\Stt S^\ABS_4}[\eta(1)] +^\ABS \eta(1)]\bigr]} \sqcup^\ABS {}$
    $\mathM{\mathS{\Stt S^\ABS_7}\bigl[\eta(0) \mapsto^\ABS \mathS{\Stt S^\ABS_7}[\eta(0)] \times^\ABS \mathS{\Stt S^\ABS_7}[\mathS{\Stt S^\ABS_7}[\eta(1)] +^\ABS \eta(1)]\bigr] \sqcup^\ABS \mathS{\Stt S^\ABS_8}}$ in
  $\mathM{\mathcal{F}^\ABS_{=0}(\mathS{v_1}, \Stt S<\ABS>) \sqcup^\ABS {}
         {\Intp MS<\ABS>}_{\Intp ST}([\mathS{\Stt S^\ABS_9}[\eta(1) \mapsto^\ABS \mathS{\Stt S^\ABS_9}[\eta(1)] +^\ABS \eta(2)]])}$
\end{metcode}
% \begin{metcode}
% fun $\mathit{cache}$ $\to$
%   fun $k$ $\to$
%     if $k = 0$
%     then
%       fun $\Stt S<\ABS>$ $\to$
%         let $\mathS{v_1^\ABS}$ =
%           $\Stt S<\ABS>[\Stt S<\ABS>[\eta(0)]]^\ABS$ in
%         let $\mathS{{\Stt S}_1^\ABS}$ = $\mathcal{F}^\ABS_{\ne 0}(\mathS{v_1^\ABS}, \Stt S<\ABS>)$ in
%         let $\mathS{v_2^\ABS}$ =
%           $\mathS{{\Stt S}_1^\ABS}[\mathS{{\Stt S}_1^\ABS}[\eta(0)]]^\ABS \mathbin{=^\ABS} \eta(1)$ in
%         let $\mathS{{\Stt S}_2^\ABS}$ = $\mathcal{F}^\ABS_{\ne 0}(\mathS{v_2^\ABS}, \mathS{{\Stt S}_1^\ABS})$ in
%         let $\mathS{{\Stt S}_3^\ABS}$ = $\mathcal{F}^\ABS_{=0}(\mathS{v_2^\ABS}, \mathS{{\Stt S}_1^\ABS})$ in
%         let $\mathS{v_3^\ABS}$ =
%           $\mathS{{\Stt S}_3^\ABS}[\mathS{{\Stt S}_3^\ABS}[\eta(0)]]^\ABS \mathbin{=^\ABS} \eta(2)$ in
%         let $\mathS{{\Stt S}_4^\ABS}$ = $\mathcal{F}^\ABS_{\ne 0}(\mathS{v_3^\ABS}, \mathS{{\Stt S}_3^\ABS})$ in
%         let $\mathS{{\Stt S}_5^\ABS}$ = $\mathcal{F}^\ABS_{=0}(\mathS{v_3^\ABS}, \mathS{{\Stt S}_3^\ABS})$ in
%         let $\mathS{{\Stt S}_6^\ABS}$ =
%           $\mathS{{\Stt S}_2^\ABS}[\eta(1) \mapsto^\ABS \mathS{{\Stt S}_2^\ABS}[\eta(1)]^\ABS \mathbin{+^\ABS}$
%                $\mathS{{\Stt S}_2^\ABS}[\mathS{{\Stt S}_2^\ABS}[\eta(0)]^\ABS \mathbin{+^\ABS} \eta(1)]^\ABS] \mathbin{\sqcup^\ABS}$
%             $(\mathS{{\Stt S}_4^\ABS}[\eta(1) \mapsto^\ABS \mathS{{\Stt S}_4^\ABS}[\eta(1)]^\ABS \mathbin{\times^\ABS}$
%                   $\mathS{{\Stt S}_4^\ABS}[\mathS{{\Stt S}_4^\ABS}[\eta(0)]^\ABS \mathbin{+^\ABS} \eta(1)]^\ABS] \mathbin{\sqcup^\ABS} \mathS{{\Stt S}_5^\ABS})$ in
%         $\mathcal{F}^\ABS_{=0}(\mathS{v_1^\ABS}, \Stt S<\ABS>) \mathbin{\sqcup^\ABS}$
%           $((\mathit{cache}\: 0)$
%              $(\mathS{{\Stt S}_6^\ABS}[\eta(0) \mapsto^\ABS \mathS{{\Stt S}_6^\ABS}[\eta(0)]^\ABS \mathbin{+^\ABS} \eta(2)]))$
%     else raise Not_found
% \end{metcode}

%\begin{metcode}
%let $\mathM{{\Intp MS<\ABS>}_{\Intp ST}(\textS{i})}$ =
%  let $\mathM{i^\ABS}$ = $\mathM{\eta(\textS{i})}$ in let $\mathM{p^\ABS}$ = $\mathM{\pi_1(\pi_1(i^\ABS)) \mathbin{=^\ABS} \eta(0)}$ in
%  $\mathM{\mathcal{F}^\ABS_{\ne 0}(p^\ABS, \pi_2(\pi_1(i^\ABS)) \mathbin{+^\ABS} \pi_2(i^\ABS)) \mathbin{\sqcup^\ABS} \mathcal{F}^\ABS_{= 0}(p^\ABS, \pi_2(\pi_1(i^\ABS)) \mathbin{\times^\ABS} \pi_2(i^\ABS)))}$
%\end{metcode}
\noindent where the subscript~$\Intp ST$ indicates that the source code~$\Intp MS<\ABS>$ is specialized with respect to the concrete interpreter~$\Intp ST$.

Notice that while the retargeted abstract interpreter~$\Intp MT<\ABS>$ is written in the meta-language~\MET, the (abstract-)interpretative overhead is eliminated and the structure of the interpreter~$\Intp*ST$ is closely mirrored:
two abstract joins in lines~9--10 mirror the if-elif-else structure of the loop body, and the recursive call in line~11 mirrors the while loop of \(\Intp*ST\).
This is a classical phenomenon in partial evaluation \citep{Jon96Introduction}.

Moreover, observe the key advantage of our approach: $\Intp MT<\ABS>$ reuses the abstract operators~$+^\ABS$ (line~9) and~$\times^\ABS$ (line~10) from the existing~$\Intp MS<\ABS>$ to analyze the \TGT*programs~\Tadd{n} and~\Tmult{n}, resulting in a sound analyzer (\hyperref[thm:correctness]{main theorem}).
Our approach enjoys the inherited properties from the existing (abstract) interpreter, as observed by \citet{Rey72Definitional}.


\section{Related Work}\label{sec:back}
\paragraph{Meta-level Analysis}
The effort to automatically keep the analysis of JavaScript (JS) programs in sync with the rapidly evolving ECMA-262~\citep{ECM15ECMAScript} specification motivated the development of a meta-level static analyzer \textsf{JSAVER}~\citep{ParAnRyu22Automatically}.
\textsf{JSAVER} targets a low-level language~\IRES, which is used to describe the semantics of JS.
By applying a meta-level analysis to a definitional interpreter of JS written in \IRES, \textsf{JSAVER} can analyze JS programs.
In this specialized setting, JS and \IRES{} share the same base value domain.

We extend this work by
\begin{enumerate}
  \item providing a language-agnostic framework, even allowing source and target languages to have different value domains; and
  \item addressing the performance limitations mentioned in their work through partial evaluation.
\end{enumerate}

\paragraph{Partial Evaluation}
It is well known that partial evaluation of an interpreter with respect to a program leads to compilation of the program---this is called the first \citet{Fut71Partial,Fut99Partial} projection.
The compiled program is essentially a specialized interpreter that does not need to traverse the AST of the original program, which eliminates interpretative overhead \citep{Jon96Introduction}.

The same idea of improving analysis performance by applying the first Futamura projection to abstract definitional interpreters \citep{DarLabNguVan17Abstracting} was demonstrated in staged abstract interpreters \citep{WeiCheRom19Staged}.
Since the analyzer is specialized to the target program, the (abstract) interpretative overhead is eliminated, resulting in efficient analysis.

\Citet{PerGalSag98Analysisa} proposed a method for partially evaluating a definitional interpreter for an imperative language written in constraint logic programming.
The resulting specialized interpreter---an imperative program compiled into a constraint logic program---can then be analyzed using an existing static analyzer for logic programming languages.
While their approach produces an efficient representation of individual target programs, our approach yields a general analyzer applicable to any program in the target language.

Note that we stage an abstract interpreter with respect to a \emph{definitional interpreter} for a new language rather than an analysis target program, representing a fundamentally different use of partial evaluation compared to prior work.

\paragraph{Reusing Abstractions}
Reusing abstractions for static analysis is an ongoing challenge.
Skeletal semantics \citep{BodGarJenSch19Skeletal} can be used to derive abstract interpreters by integrating meta-language abstractions with language-specific ones \citep{JenRebSch23Deriving}.
Composing modular analysis components has been proposed by \citet{KeiErd19Sound}, which employs arrows meta-language \citep{Hug00Generalising}.
An open-source static analysis framework \textsc{Mopsa} \citep{JouMinMonOua20Combinations} leverages OCaml's extensible variants to easily define new targets and reuse abstractions compositionally.

Our method follows a top-down approach of recycling an entire abstract interpreter to retarget to a new one, instead of composing analysis components in a bottom-up fashion.


\section{Discussion}\label{sec:discuss}
Our recipe for retargeting abstract interpreters can drastically reduce the effort required to develop static analyzers for various applications.
For example, a JavaScript static analyzer such as JSAVER~\citep{ParAnRyu22Automatically} can be retargeted to analyze React applications using a React definitional interpreter like \textsc{\textsf{React-tRace}}~\citep{LeeAhnYi25React}.

We are currently working on identifying the properties of the retargeted abstract interpreter---for example, the effect of the semantic gap between the source and target languages, its precision and performance characteristics, and approaches to inheriting analysis sensitivity for the target language.
We believe that our recipe can be extended to provide a knob for analysis designers, enabling sensitivity control over the analyses of target programs.


\begin{acks}
We thank John Patrick Gallagher and anonymous reviewers for their valuable feedback.

A preliminary version of this work was presented at the Student Research Competition (SRC) at the 46th ACM SIGPLAN Conference on Programming Language Design and Implementation (PLDI~'25) \citep{Lee25Retargeting}.

This work was supported by  the National Research Foundation of Korea (NRF) grant funded by the Korea government (MSIT) (0536-20250107), BK21 FOUR Intelligence Computing (Dept. of Computer Science and Engineering, SNU) funded by National Research Foundation of Korea (NRF) (4199990214639), Greenlabs Co., Ltd. (0536-20220078), and Samsung Electronics Co., Ltd. (0536-20230088).
\end{acks}

\bibliographystyle{ACM-Reference-Format}
\bibliography{references}
\end{document}
